% !TEX program = xelatex
% =============================================================================
%  CUMCM XeLaTeX 模板 —— 可编译演示样例（standalone demo）
%
%  本文件独立于「正式模板」skills/5writing/templates/zh/cumcm-latex/，
%  不会随模板实例化进入每个 run。它的作用是：
%    1. 演示模板的排版效果（公式 / 三线表 / 跨页长表 / 图 / 代码 / 交叉引用 / 参考文献 / 附录）；
%    2. 作为可编译的回归样例，验证模板 preamble 在本机 XeLaTeX 下确实通过。
%
%  preamble 与正式模板 main.tex 保持一致（仅额外加载 tikz 用于自绘演示图，
%  避免演示依赖外部图片）。修改模板 preamble 时，请同步本文件。
%
%  编译：xelatex main_demo.tex（交叉引用需两遍）。
% =============================================================================
\documentclass[UTF8,fontset=none,12pt,a4paper]{ctexart}

% ============ 论文配置块（演示值） ============
\newcommand{\PaperTitle}{基于相位—波数关系的外延层厚度红外反演模型（演示样例）}
\newif\ifpaperhasreferences   \paperhasreferencestrue
\newif\ifnumberedreferences   \numberedreferencesfalse   % 改 true 可演示“编号参考文献”
\newif\ifshowappendix         \showappendixtrue

% ---------- 字体：Windows / macOS / Linux 自动回退，最终兜底 Fandol ----------
\usepackage{fontspec}
\usepackage{xeCJK}
\def\zhSong{FandolSong}\def\zhHei{FandolHei}\def\zhKai{FandolKai}\def\zhFang{FandolFang}
\IfFontExistsTF{Noto Serif CJK SC}{\def\zhSong{Noto Serif CJK SC}}{}
\IfFontExistsTF{Noto Sans CJK SC}{\def\zhHei{Noto Sans CJK SC}}{}
\IfFontExistsTF{Songti SC}{\def\zhSong{Songti SC}}{}
\IfFontExistsTF{Heiti SC}{\def\zhHei{Heiti SC}}{}
\IfFontExistsTF{Kaiti SC}{\def\zhKai{Kaiti SC}}{}
\IfFontExistsTF{STFangsong}{\def\zhFang{STFangsong}}{}
\IfFontExistsTF{SimSun}{\def\zhSong{SimSun}}{}
\IfFontExistsTF{SimHei}{\def\zhHei{SimHei}}{}
\IfFontExistsTF{KaiTi}{\def\zhKai{KaiTi}}{}
\IfFontExistsTF{FangSong}{\def\zhFang{FangSong}}{}
\setCJKmainfont{\zhSong}
\setCJKsansfont{\zhHei}
\setCJKmonofont{\zhFang}
\newCJKfontfamily\paperSong{\zhSong}
\newCJKfontfamily\paperHei{\zhHei}
\newCJKfontfamily\paperKai{\zhKai}
\newCJKfontfamily\paperFang{\zhFang}
\let\songti\paperSong
\let\heiti\paperHei
\let\kaishu\paperKai
\let\fangsong\paperFang
\IfFontExistsTF{Times New Roman}{\setmainfont{Times New Roman}}{%
  \IfFontExistsTF{Liberation Serif}{\setmainfont{Liberation Serif}}{}%
}
\IfFontExistsTF{Consolas}{\setmonofont{Consolas}}{%
  \IfFontExistsTF{Liberation Mono}{\setmonofont{Liberation Mono}}{}%
}

% ---------- 页面与正文 ----------
\usepackage[a4paper, top=2.5cm, bottom=2.5cm, left=2.5cm, right=2.5cm, footskip=1.2cm]{geometry}
\usepackage{setspace}
\setstretch{1.40}
\setlength{\parindent}{2em}
\setlength{\parskip}{0pt}
\raggedbottom

\usepackage{fancyhdr}
\pagestyle{fancy}
\fancyhf{}
\fancyfoot[C]{\fontsize{10.5pt}{12pt}\selectfont\thepage}
\renewcommand{\headrulewidth}{0pt}
\renewcommand{\footrulewidth}{0pt}

\usepackage{amsmath,amssymb,mathtools,bm}
\allowdisplaybreaks[2]

\usepackage{graphicx}
\graphicspath{{figures/}{figures/current/}{../figures/current/}}
\usepackage{booktabs,tabularx,longtable,array,multirow}
\usepackage{float}
\usepackage{placeins}
\usepackage{caption}
\DeclareCaptionFont{cncaption}{\fontsize{10.5pt}{14pt}\selectfont}
\captionsetup{font=cncaption,labelfont=bf,labelsep=space,justification=centering,singlelinecheck=true}
\captionsetup[figure]{position=bottom,skip=6pt}
\captionsetup[table]{position=top,skip=6pt}
\renewcommand{\arraystretch}{1.20}
\setlength{\tabcolsep}{6pt}

\usepackage{enumitem}
\setlist{nosep,leftmargin=2em}
\setlist[enumerate,1]{label=\arabic*、}

\ctexset{
  section = {name={,、},number=\chinese{section},format=\centering\heiti\fontsize{15pt}{20pt}\selectfont,aftername=\hspace{0.5em},beforeskip=18pt,afterskip=10pt,afterindent=false},
  subsection = {name={,},number=\arabic{section}.\arabic{subsection},format=\heiti\fontsize{12pt}{18pt}\selectfont,aftername=\hspace{0.5em},beforeskip=10pt,afterskip=2pt,afterindent=false},
  subsubsection = {name={,},number=\arabic{section}.\arabic{subsection}.\arabic{subsubsection},format=\heiti\fontsize{12pt}{18pt}\selectfont,aftername=\hspace{0.5em},beforeskip=8pt,afterskip=2pt,afterindent=false}
}

\usepackage{xcolor}
\usepackage{listings}
\lstset{
  basicstyle=\ttfamily\fontsize{9pt}{11pt}\selectfont,
  keywordstyle=\bfseries,commentstyle=\itshape,
  numbers=left,numberstyle=\tiny,numbersep=8pt,
  frame=single,rulecolor=\color{black!40},framerule=0.4pt,framesep=5pt,
  breaklines=true,breakatwhitespace=false,showstringspaces=false,
  columns=fullflexible,keepspaces=true,tabsize=4,aboveskip=8pt,belowskip=8pt
}

\usepackage[hidelinks]{hyperref}

% ---------- 参考文献标题单一来源（与模板一致） ----------
\usepackage{etoolbox}
\AtBeginDocument{%
  \ifnumberedreferences
    \patchcmd{\thebibliography}{\section*{\refname}}{\section{\refname}}{}{%
      \PackageWarning{cumcm-template}{无法为参考文献标题加编号：thebibliography 结构与预期不符}%
    }%
  \fi
}

% ---------- 仅演示需要：自绘示意图，避免依赖外部图片 ----------
\usepackage{tikz}
\usetikzlibrary{positioning,arrows.meta}

% ---------- 自定义命令（与模板一致） ----------
\newcommand{\papertitle}[1]{%
  \vspace*{0.55cm}%
  \begin{center}{\heiti\fontsize{18pt}{24pt}\selectfont #1\par}\end{center}
  \vspace{0.65cm}%
}
\newcommand{\abstractcn}[2]{%
  \begin{center}{\heiti\fontsize{15pt}{20pt}\selectfont 摘要\par}\end{center}
  \vspace{0.35em}%
  \begingroup
    \setstretch{1.32}\fontsize{12pt}{18pt}\selectfont
    #1\par
    \vspace{0.65em}%
    \noindent{\heiti 关键词：}#2\par
  \endgroup
  \clearpage
}
\newcommand{\threelinetable}[5][]{%
  \begin{table}[htbp]
    \centering
    \caption{#2}%
    \if\relax\detokenize{#1}\relax\else\label{#1}\fi
    \begin{tabular}{#3}\toprule #4 \\ \midrule #5 \\ \bottomrule\end{tabular}
  \end{table}
}
\newcounter{paperappendix}
\newcommand{\appendixsection}[1]{\stepcounter{paperappendix}\phantomsection\section*{附录~\Alph{paperappendix}\quad #1}}
\newcommand{\referencescn}{%
  \ifpaperhasreferences
    \FloatBarrier
    \IfFileExists{references.tex}{%
      \begingroup\setstretch{1.18}\fontsize{10.5pt}{14pt}\selectfont% 演示参考文献。用 thebibliography 环境；标题“参考文献”由环境自动生成，
% 请勿再手写“参考文献”标题（否则会出现两个标题）。
\begin{thebibliography}{99}
\bibitem{ref-fowles} Fowles G R. Introduction to Modern Optics[M]. 2nd ed. New York: Dover, 1989.
\bibitem{ref-born} Born M, Wolf E. Principles of Optics[M]. 7th ed. Cambridge: Cambridge University Press, 1999.
\bibitem{ref-harrick} Harrick N J. Internal Reflection Spectroscopy[M]. New York: Interscience, 1967.
\end{thebibliography}
\endgroup
    }{\PackageError{cumcm-template}{references.tex is required when paperhasreferences is true}{Provide references.tex, or set \noexpand\paperhasreferencesfalse.}}%
  \fi
}
\newcommand{\appendixcn}{%
  \FloatBarrier
  \IfFileExists{sections/A_code.tex}{\clearpage\appendixsection{核心代码}% 演示附录代码。真实论文只放核心片段，完整代码作为支撑材料另附。
\begin{lstlisting}[language=Python]
import numpy as np
from scipy.optimize import least_squares

def phase_order_fit(nu_peaks, n_of_nu, theta_deg):
    """双角度共享厚度的峰序号回归（演示）。"""
    theta = np.deg2rad(theta_deg)
    def g(nu):
        n = n_of_nu(nu)
        return n * nu * np.sqrt(1.0 - (np.sin(theta) / n) ** 2)
    m = np.arange(len(nu_peaks))            # 峰序号
    gvals = g(nu_peaks)
    # m = a + 2 d g(nu)  -> 线性最小二乘求 d
    A = np.vstack([np.ones_like(gvals), 2.0 * gvals]).T
    (a, d), *_ = np.linalg.lstsq(A, m, rcond=None)
    return d
\end{lstlisting}
}{}%
}

\begin{document}
\setcounter{page}{1}

\papertitle{\PaperTitle}

\abstractcn{%
  红外干涉测厚的核心困难在于：反射谱同时受层厚、折射率色散与多光束效应影响，
  三者在谱形上高度耦合，直接读峰会把色散误差混入厚度。本文将该问题定义为
  “以相位—波数关系为约束的联合反演”，以双角度共享厚度的峰序号回归为主方法，
  配合分段一致性检验与残差块自助法量化不确定度。在演示数据上得到外延层厚度
  点估计 $8.11\ \mu\mathrm{m}$，$95\%$ 置信区间约 $[7.77,\,8.25]\ \mu\mathrm{m}$，
  重构 RMSE 约 $1.5\%$。分析表明：在理想无吸收单层假设下，多光束改变谱形与对比度，
  但不改变同类极值的相位间距，故峰间距法的解析厚度偏差为零；结论在折射率
  $\pm3\%$ 扰动内稳定，当界面平行性假设被打破时需重新评估。%
}{%
  红外干涉\quad 相位—波数关系\quad 峰序号回归\quad 自助法（Bootstrap）%
}

% 演示正文：展示公式、三线表、跨页长表、示意图、代码、交叉引用等模板能力。

\section{问题重述与分析}
本演示以红外干涉测厚为例。给定不同入射角下的反射谱，需反演外延层厚度 $d$，
并评估折射率色散与多光束效应对结果的影响。核心变量为厚度 $d$、折射率 $n(\nu)$
与波数 $\nu$；难点在于三者在谱形上的耦合。整体求解路线见图~\ref{fig:roadmap}。

\begin{figure}[htbp]
  \centering
  \begin{tikzpicture}[
    node distance=8mm,
    box/.style={draw, rounded corners, align=center, inner sep=4pt, font=\small},
    arr/.style={-latex, thick}]
    \node[box] (a) {反射谱\\预处理};
    \node[box, right=of a] (b) {峰检测与\\级次分配};
    \node[box, right=of b] (c) {双角度\\联合回归};
    \node[box, right=of c] (d) {厚度 $d$\\与不确定度};
    \draw[arr] (a) -- (b);
    \draw[arr] (b) -- (c);
    \draw[arr] (c) -- (d);
  \end{tikzpicture}
  \caption{总体求解路线示意图}
  \label{fig:roadmap}
\end{figure}

\section{模型的建立与求解}

\subsection{相位—波数关系}
设界面反射系数为 $r_1,r_2$，相位差 $\delta(\nu)=4\pi d\, n(\nu)\nu\cos\theta'$。
双光束近似下反射率为
\begin{equation}
  R_{2b}(\nu)=A(\nu)+B(\nu)\cos\delta(\nu),
  \label{eq:twobeam}
\end{equation}
其中 $A,B$ 为慢变包络。式~\eqref{eq:twobeam} 中相邻同类极值满足相位差 $2\pi$，
据此可由峰序号回归稳健地估计 $d$。作为对照，完整 Airy 反射率为
\begin{equation}
  R_{A}(\nu)=\frac{r_1^2+r_2^2+2r_1r_2\cos\delta}{1+r_1^2r_2^2+2r_1r_2\cos\delta}.
  \label{eq:airy}
\end{equation}
式~\eqref{eq:airy} 的分母正是多次反射求和的结果，与式~\eqref{eq:twobeam} 的双光束近似
在物理含义上必须区分，不能混用。

\subsection{结果与检验}
演示数据上的主结果汇总于表~\ref{tab:main}。可见两角度点估计一致，
联合区间较单角度显著收窄。

\threelinetable[tab:main]{主结果汇总（演示数据）}{lccc}
{方法 & 角度 & 点估计 $d/\mu\mathrm{m}$ & 95\% CI $/\mu\mathrm{m}$}
{峰序号回归 & $10^\circ$ & 8.09 & $[7.36,\,8.53]$ \\
 峰序号回归 & $15^\circ$ & 8.14 & $[7.33,\,9.38]$ \\
 双角联合 & 共享 & 8.11 & $[7.77,\,8.25]$}

为检验分段一致性，将全谱划分为若干波段分别估计厚度，结果见表~\ref{tab:seg}
（长表演示，可跨页）。

\begin{longtable}{cccc}
  \caption{分段一致性检验（演示数据，长表跨页示例）}\label{tab:seg}\\
  \toprule
  波段编号 & 波段中心 $/\mathrm{cm^{-1}}$ & 厚度 $d/\mu\mathrm{m}$ & 相对偏差 \\
  \midrule
  \endfirsthead
  \multicolumn{4}{l}{\small（续表~\ref{tab:seg}）}\\
  \toprule
  波段编号 & 波段中心 $/\mathrm{cm^{-1}}$ & 厚度 $d/\mu\mathrm{m}$ & 相对偏差 \\
  \midrule
  \endhead
  \midrule
  \multicolumn{4}{r}{\small 接下页}\\
  \endfoot
  \bottomrule
  \endlastfoot
  1 & 1050 & 8.02 & $-1.1\%$ \\
  2 & 1150 & 8.09 & $-0.2\%$ \\
  3 & 1250 & 8.13 & $+0.2\%$ \\
  4 & 1350 & 8.15 & $+0.5\%$ \\
  5 & 1450 & 8.10 & $-0.1\%$ \\
  6 & 1550 & 8.07 & $-0.5\%$ \\
  7 & 1650 & 8.12 & $+0.1\%$ \\
  8 & 1750 & 8.14 & $+0.4\%$ \\
\end{longtable}

核心求解步骤的代码见附录~A。

\section{模型评价、改进与推广}
本模型的优点是以相位—波数关系为统一约束，避免了逐峰读数对色散的敏感；
双角度联合回归在不牺牲物理可解释性的前提下收窄了区间。局限在于假设界面平行、
无吸收，若样品存在明显吸收带或界面倾斜，需引入 $A(\nu),B(\nu)$ 的更高阶描述。
该框架可推广至多层结构的分层反演，以及其它基于干涉条纹的测厚场景。


\referencescn

\ifshowappendix
  \appendixcn
\fi

\end{document}
